% Part: lambda-calculus
% Chapter: lambda-definability
% Section: lists

\documentclass[../../../include/open-logic-section]{subfiles}

\begin{document}

\olfileid{lam}{ldf}{lst}
\olsection{列表}

列表 $[M_1, M_2, \ldots, M_n]$ 可定义为：
\[
  [M_1, M_2, \ldots, M_n] \ident \lambd[sz][s M_1 (s M_2 (\ldots (s M_n z)))]
\]
非形式地说，表示列表的项就是该列表的“累加器”函数：它接受两个项作为自变元；第一个自变元是一个函数，该函数接收当前的累加值和新元素，并返回新的累加值；第二个自变元是初始的累加值。该函数逐个累加元素，最终返回结果。

\begin{digress}
  如果你熟悉函数式编程，会发现此定义类似于\emph{右折叠}：列表被定义为其所含元素的右折叠函数。
\end{digress}

基于这种编码，我们很容易定义一些有用的函数；
\begin{align*}
  \fn{Sum} & \ident
  \lambd[l][l \, (\lambd[xa][\fn{Add} \, x \, a])\,  \num{0}]\\
  \fn{Len} & \ident
  \lambd[l][l \, (\lambd[xa][\fn{Add} \, a \, \num{1}]) \num{0}]
\end{align*}

$\fn{Sum}$ 计算邱奇数列表的和。它的工作方式是对列表进行累加，其中初始值为 $\num{0}$；对每个元素，将 $\fn{Add} \, x\, a$ 作为新值返回（$x$ 为当前元素，$a$ 为累加值）。最终结果即为所有元素的总和。

\begin{prob}
  $\fn{Len}$ 的作用是什么？请解释。
\end{prob}

\begin{prob}
  定义一个反转列表的函数。
\end{prob}
\end{document}